\chapter{Coding Theory and the Hamming Code}
\label{ch:coding-theory-hamming}

\section{Optional Project}

This chapter is a compact project. It shows how kernels, matrix-vector
products, and subspaces work over the two-element field
\[
  \Z_2=\{0,1\},
\]
where addition and multiplication are done modulo $2$.

\section{Messages, Codewords, And Errors}

A 4-bit message is a vector in $\Z_2^4$. A 7-bit transmitted word is a vector
in $\Z_2^7$. The goal of the Hamming$(7,4)$ code is to encode 4 data bits into
7 bits so that any single flipped bit can be detected and corrected.

The extra 3 bits are redundancy. Linear algebra organizes that redundancy.

\section{The Parity-Check Matrix}

Define
\[
  H=
  \begin{bmatrix}
  0&0&0&1&1&1&1\\
  0&1&1&0&0&1&1\\
  1&0&1&0&1&0&1
  \end{bmatrix}.
\]
The columns are the nonzero binary vectors in $\Z_2^3$:
\[
  001,\ 010,\ 011,\ 100,\ 101,\ 110,\ 111.
\]

\begin{definition}[Code]
The Hamming$(7,4)$ code is
\[
  \mathcal{C}=\Ker(H)
  =
  \{c\in\Z_2^7: Hc=0\}.
\]
Vectors in $\mathcal{C}$ are called codewords.
\end{definition}

Since $H$ has rank $3$, rank-nullity gives
\[
  \dim\Ker(H)=7-3=4.
\]
So $\mathcal{C}$ has $2^4=16$ codewords, one for each 4-bit message.

\section{A Generator Matrix}

A generator matrix has columns forming a basis for $\Ker(H)$. One choice is
\[
  G=
  \begin{bmatrix}
  1&0&0&0\\
  0&1&0&0\\
  0&0&1&0\\
  0&0&0&1\\
  1&1&0&1\\
  1&0&1&1\\
  0&1&1&1
  \end{bmatrix}.
\]
Encoding a message $m\in\Z_2^4$ means computing
\[
  c=Gm\pmod 2.
\]
The first four bits of $c$ are the original message. The last three bits are
parity bits.

\begin{proposition}
Every encoded vector $c=Gm$ is a codeword.
\end{proposition}

\begin{proof}
It is enough to check
\[
  HG=0\pmod 2.
\]
Then
\[
  Hc=H(Gm)=(HG)m=0.
\]
\end{proof}

\section{Syndromes}

Suppose a codeword $c$ is transmitted and the received vector is
\[
  r=c+e,
\]
where $e$ is an error vector. Compute
\[
  s=Hr.
\]
This vector $s\in\Z_2^3$ is called the \emph{syndrome}.

Since $Hc=0$,
\[
  s=H(c+e)=Hc+He=He.
\]
Thus the syndrome depends only on the error.

\begin{theorem}[Single-bit correction]
If at most one bit is flipped, then the syndrome either equals $0$ or equals
the column of $H$ corresponding to the flipped bit.
\end{theorem}

\begin{proof}
If there is no error, then $e=0$ and $s=He=0$. If bit $k$ is flipped, then
$e=e_k$, the $k$-th standard basis vector. Therefore
\[
  s=He_k,
\]
which is the $k$-th column of $H$. The columns of $H$ are all nonzero and
distinct, so the syndrome identifies the flipped bit.
\end{proof}

\begin{warningbox}
This code corrects one flipped bit. Two flipped bits can produce a misleading
syndrome, because the syndrome is the sum of two columns. The guarantee is
single-error correction, not arbitrary error correction.
\end{warningbox}

\section{Python Thread}

\begin{lstlisting}
import numpy as np

def mod2(x):
    return x % 2

H = np.array([[0,0,0,1,1,1,1],
              [0,1,1,0,0,1,1],
              [1,0,1,0,1,0,1]], dtype=int)

G = np.array([[1,0,0,0],
              [0,1,0,0],
              [0,0,1,0],
              [0,0,0,1],
              [1,1,0,1],
              [1,0,1,1],
              [0,1,1,1]], dtype=int)

print("H @ G mod 2:")
print(mod2(H @ G))

m = np.array([1, 0, 1, 1], dtype=int)
c = mod2(G @ m)
print("message:", m)
print("codeword:", c)

# Flip bit 5 in one-indexed language, index 4 in Python.
r = c.copy()
r[4] = 1 - r[4]

s = mod2(H @ r)
print("received:", r)
print("syndrome:", s)

# Build a syndrome lookup table from columns of H.
lookup = {tuple(H[:, j]): j for j in range(7)}
error_index = lookup[tuple(s)]

fixed = r.copy()
fixed[error_index] = 1 - fixed[error_index]
print("corrected:", fixed)
print("matches:", np.array_equal(fixed, c))
\end{lstlisting}

\section{What This Project Reinforces}

\begin{itemize}
  \item A code can be a null space.
  \item Redundancy can be organized by a parity-check matrix.
  \item Matrix-vector multiplication can detect violations of parity rules.
  \item A syndrome is an image vector that identifies a likely error.
  \item The same rank-nullity theorem works over $\Z_2$.
\end{itemize}

\section*{Chapter Summary}
\addcontentsline{toc}{section}{Chapter Summary}

\begin{itemize}
  \item Hamming$(7,4)$ encodes 4-bit messages as 7-bit codewords.
  \item The code is $\Ker(H)$ over $\Z_2$.
  \item The generator matrix maps messages into the code.
  \item The syndrome $Hr$ identifies a single flipped bit.
  \item This is a compact application of kernels, images, and matrix products.
\end{itemize}

\section*{Exercises}
\addcontentsline{toc}{section}{Exercises}
\subsection*{A. Theory}
\begin{exercise}\exerciseid{17.A1} State the Eckart-Young theorem in your own words. \end{exercise}
\subsection*{B. By Hand}
\begin{exercise}\exerciseid{17.B1} If the singular values of $A$ are $5,3,1$, compute the Frobenius error of the best rank-$1$ approximation. \end{exercise}
\subsection*{C. Python / NumPy}
\begin{exercise}\exerciseid{17.C1} Compute rank-$k$ approximations of a matrix using \texttt{np.linalg.svd}. Plot or print the errors. \end{exercise}
\subsection*{D. Challenge}
\begin{exercise}\exerciseid{17.D1} Explain why PCA begins by centering the data matrix. \end{exercise}

